\section{Experience}

% Each header row is one line of text, so parsers read the stack and dates with the role.
\begin{onecolentry}
    \textbf{Software Engineer} \textbar\ \textit{Callab AI (YC P26)}\hfill\textit{Tunis, Tunisia (On Site)}

    {\footnotesize TypeScript, Python, React, Bun, AI SDK, MCP, PostgreSQL, Redis, Docker, Kubernetes, AWS}\hfill\textit{Jun 2025 – Present}%
\end{onecolentry}
\vspace{0.03 cm}
\begin{onecolentry}
    \begin{highlights}
        \item \textbf{Shipped} features \textbf{end to end} across backend, AI, frontend and DevOps, and \textbf{owned production systems} in a multi-service platform, delivering dozens of features and resolving production incidents and bugs.
        \item \textbf{Built} the \textbf{Omnichannel AI Agent} serving Slack, Telegram, Teams, WhatsApp and web, with \textbf{human-in-the-loop approvals}, a flow engine, \textbf{long-term memory}, scheduled tasks, knowledge-base retrieval, MCP tools and skills, and conversation compaction.
        \item \textbf{Delivered human call supervision} end to end, letting operators watch live calls, listen in and \textbf{take over} from the \textbf{AI agent} mid-conversation, with agents escalating to a human when they cannot resolve a call.
        \item \textbf{Built} an \textbf{in-product AI assistant} that answers product questions and drives the app for the user, \textbf{reading and acting} on the \textbf{live DOM} through a browser tool family, so it can navigate, fill and complete tasks on the user's behalf.
    \end{highlights}
\end{onecolentry}
\entrygap

\begin{onecolentry}
    \textbf{Google Summer of Code Developer} \textbar\ \textit{NRNB (National Resource for Network Biology)}\hfill\textit{Remote}

    {\footnotesize Python, TypeScript, FastAPI, Next.js, Qdrant, Langfuse}\hfill\textit{May 2025 – Sep 2025}%
\end{onecolentry}
\vspace{0.03 cm}
\begin{onecolentry}
    \begin{highlights}
        \item Collaborated with the VCell team to \textbf{develop and deploy \href{https://github.com/virtualcell/VCell-AI}{VCell-AI}}, an \textbf{AI agent platform} with tool calling capabilities, built on a Next.js frontend and FastAPI backend, enabling researchers to \textbf{query, explore, and create biological models}.
        \item \textbf{Integrated} a \textbf{RAG} system using \textbf{Qdrant} to provide contextual information about VCell software, tutorials, and guidelines, with \textbf{LLM observability} via \textbf{Langfuse} to track agent interactions, monitor performance, and improve reliability.
        \item \textbf{Enabled} the AI agent to \textbf{analyze diagrams and VCML files} using \textbf{multimodal} capabilities, supporting bio-model interpretation.
    \end{highlights}
\end{onecolentry}
\entrygap

\begin{onecolentry}
    \textbf{Machine Learning Engineering Intern} \textbar\ \textit{Orange}\hfill\textit{Tunis, Tunisia (On Site)}

    {\footnotesize Python, PyTorch, CrewAI, Streamlit, Flask}\hfill\textit{Aug 2024 – Oct 2024}%
\end{onecolentry}
\vspace{0.03 cm}
\begin{onecolentry}
    \begin{highlights}
        \item \textbf{Developed} a lightweight \textbf{CNN model} for plant disease classification optimized for \textbf{edge devices} using \textbf{knowledge distillation}.
        \item \textbf{Built} \textbf{AI agentic workflow} for \textbf{automated report generation}, streamlining data collection, analysis, and research extraction.
        \item \textbf{Achieved 4th place} in \textbf{Orange Summer Challenge}, pitching the solution to both technical and non-technical stakeholders.
    \end{highlights}
\end{onecolentry}
\vspace{0.1 cm}
